\scanpage{191}
\margin{\unsure{1}}
\begin{center}- \mspage{62} -\end{center}
\margin{IX}
\margin{$=$}
\margin{-}

où on se rappellera que les groupes $T_\ell(A_K)^f$ et $T_\ell(A'_K)^f$ qui figurent au second membre sont les restrictions à $\eta$ de pro-$\ell$-groupes de Barsotti-Tate définis sur $S$, $T_\ell(A^{\unsure{0|o}})^f$ et $T_\ell(A'^{\unsure{0|o}})^f$ respectivement ; la lettre $D$ désigne ici les duals au sens des pro-groupes de Barsotti-Tate, savoir $\operatorname{\uline{Hom}}(-,\mathbb Z_\ell(1))$. Ainsi :

\uline{Proposition} 5.6. \uline{Supposons} $S$ \uline{complet et que} $A_K$ \uline{ait réduction semi-stable sur} $S$. \uline{Alors pour tout nombre premier} $\ell$ (\uline{qui peut être égal à la caractéristique résiduelle de} $S$), \uline{le pro-}$\ell$\uline{-groupe de Barsotti-Tate} $T_\ell(A_K)^t$ \uline{est la fibre générique d'un pro-groupe de Barsotti-Tate sur} $S$, ``\uline{de type torique}'', \uline{donné par} (5.5.7), \uline{tandis que} $T_\ell(A_K)/T_\ell(A_K)^f$ \uline{est la fibre générique d'un pro-groupe de Barsotti-Tate sur} $S$ \uline{de type pro-étale}, \uline{donné par} (5.5.8). \uline{Enfin les groupes} $T_\ell(A_K)^f$ \uline{et} $T_\ell(A_K)/T_\ell(A_K)^t$ \uline{sont les fibres génériques de pro-groupes de Barsotti-Tate définis sur} $S$, \uline{égaux respectivement à} $T_\ell(A^{\unsure{0|o}})^f$ \uline{et à} $D(T_\ell(A'^{\unsure{0|o}})^f)\unsure{-|=}\operatorname{\uline{Hom}}(T_\ell(A'^{\unsure{0|o}})^f,\mathbb Z_\ell(1))$.

5.6.1. On notera qu'il revient au même pour le revêtement étale $S'$ de déployer $T_\ell(A_K)^t$ ou de déployer $T_\ell(A_K)/T_\ell(A_K)^f$, car cela signifie respectivement que $S'$ déploie les groupes constants tordus $\underline M$ et \uline{$M'$}, or ceux-ci sont isogènes, puisque $T_{\unsure{o|0}}$ et $T'_{\unsure{o|0}}$ le sont. On voit en même temps qu'il existe une plus petite extension étale $S'$ de $S$ qui fait l'affaire.

Signalons aussi la conséquence facile suivante des définitions de 2.2.3 :

\margin{380}
\margin{-}
\margin{380}
\margin{190}
